% answers_short.tex —— 2015 普通高等学校招生考试（山东卷理） · 简短答案速查
% 由 tools/make_short_answers.py 自 answers.tex 的 \gksolhead/\gkanswer 提取，
% 并与 2015/2015_数学真题_山东理_答案要点.md 逐题比对后生成，勿手改（改源 MD 后重跑）。

\begin{gknotice}
\textbf{使用说明}：本节只给最终答案便于快速核对，完整推导见解析版。答案为编写者独立推导，\textbf{非官方标准答案}。
\end{gknotice}

\gkpart{选择题}
\begin{center}\begin{tabular}{ccccccccccc}
\toprule
题号 & 1 & 2 & 3 & 4 & 5 & 6 & 7 & 8 & 9 & 10 \\
\midrule
答案 & C & A & B & D & A & B & C & B & D & C \\
\bottomrule
\end{tabular}\end{center}

\gkpart{填空题}
\begin{itemize}[leftmargin=2.6em]
  \item[\textbf{11.}] $4^{n-1}$
  \item[\textbf{12.}] $1$
  \item[\textbf{13.}] $\dfrac{11}{6}$
  \item[\textbf{14.}] $-\dfrac{3}{2}$
  \item[\textbf{15.}] $\dfrac{3}{2}$
\end{itemize}

\gkpart{解答题}
\begin{itemize}[leftmargin=2.6em]
  \item[\textbf{16.}] (1) 递增区间为 $[-\dfrac{\pi}{4}+k\pi,\dfrac{\pi}{4}+k\pi]$，递减区间为 $[\dfrac{\pi}{4}+k\pi,\dfrac{3\pi}{4}+k\pi]$，$k\in\mathbb{Z}$；(2) $\dfrac{2+\sqrt{3}}{4}$
  \item[\textbf{17.}] (1) $BD\parallel$ 平面 $FGH$；(2) $60^\circ$
  \item[\textbf{18.}] (1) $a_1=3$，$a_n=3^{n-1}\ (n\ge2)$；(2) $T_n=\dfrac{13}{12}-\dfrac{2n+1}{4\cdot3^{n-1}}$
  \item[\textbf{19.}] (1) $125,135,145,235,245,345$；(2) $P(X=-1)=\dfrac{1}{14}$，$P(X=0)=\dfrac{2}{3}$，$P(X=1)=\dfrac{11}{42}$；$EX=\dfrac{4}{21}$
  \item[\textbf{20.}] (1) $\dfrac{x^2}{4}+y^2=1$；(2) ① $\dfrac{|OQ|}{|OP|}=2$；② $\triangle ABQ$ 面积的最大值为 $6\sqrt{3}$
  \item[\textbf{21.}] (1) 当 $a<0$ 时有 1 个极值点；当 $0\le a\le\dfrac{8}{9}$ 时无极值点；当 $a>\dfrac{8}{9}$ 时有 2 个极值点；(2) $0\le a\le1$
\end{itemize}
